% Generated draft: requires mathematical and source review.
\documentclass[11pt]{article}
\usepackage[margin=25mm]{geometry}
\usepackage{fontspec}
\setmainfont{Latin Modern Roman}
\usepackage{amsmath,amssymb,mathtools,unicode-math}
\setmathfont{Latin Modern Math}
\usepackage{enumitem,needspace,xurl,hyperref}
\setlength{\parindent}{0pt}
\setlength{\parskip}{5pt}
\setlength{\emergencystretch}{3em}
\newcommand{\N}{\mathbb{N}}
\newcommand{\Z}{\mathbb{Z}}
\newcommand{\Q}{\mathbb{Q}}
\newcommand{\R}{\mathbb{R}}
\newcommand{\C}{\mathbb{C}}
\newcommand{\F}{\mathbb{F}}
\newcommand{\A}{\mathbb{A}}
\newcommand{\PP}{\mathbb{P}}
\newcommand{\E}{\mathbb{E}}
\newcommand{\eps}{\varepsilon}
\newcommand{\dd}{\,\mathrm{d}}
\begin{document}
\section*{UnsolvedMath ID: 30005174}
\textbf{Source problem code:} catalogue entry
\subsection{Definitions and mathematical statement}
Let \(X_t\) denote the smooth fibres of an \(n\)-dimensional Calabi--Yau degeneration. The supplied excerpt does not specify the base field, the parameter domain, the precise Calabi--Yau hypotheses, or the degeneration assumptions. The formulation below is therefore conditional on supplying these data and the constructions referenced in the omitted Conjecture 1.

Let \(B\) denote the limiting base of the Strominger--Yau--Zaslow (SYZ) torus fibrations referred to in that conjecture, and let \(\Delta\subset B\) be its singular locus. Write \(X\) for the associated non-Archimedean generic fibre, \(X^{\mathrm{an}}\) for its analytification, and \(\mathrm{Sk}(X)\subset X^{\mathrm{an}}\) for its essential skeleton. In the usual setting where the construction is defined, the essential skeleton is obtained from the minimum loci of weight functions associated with nonzero pluricanonical forms; the precise setting and conventions are not supplied here. For the simplicial structure used in the quoted result, put
\[
Z=\bigcup_{\substack{\sigma\text{ a simplex of }\mathrm{Sk}(X)\\ \operatorname{codim}\sigma\geq 2}}\sigma,
\qquad
B^{\mathrm{reg}}=B\setminus\Delta,
\qquad
\mathrm{Sk}(X)^{\mathrm{reg}}=\mathrm{Sk}(X)\setminus Z.
\]
An integral-affine structure on an \(n\)-manifold is an atlas whose transition maps have the form \(x\mapsto Ax+b\), with \(A\in\mathrm{GL}(n,\Z)\) and \(b\in\R^n\). An integral-affine isomorphism is a homeomorphism that is locally of this form in the respective charts. A metric integral-affine comparison must additionally preserve the specified metric structures; an integral-affine identification alone does not establish this additional requirement.

\textbf{Kontsevich--Soibelman comparison problem, as reported in the supplied record.}
For each Calabi--Yau degeneration satisfying the intended, but incompletely supplied, hypotheses, do the SYZ base and the essential skeleton agree as metric integral-affine manifolds away from their respective singular loci? In the notation of the excerpt, the intended regular-locus comparison asks for an integral-affine isomorphism
\[
\Phi:B\setminus\Delta\longrightarrow\mathrm{Sk}(X)\setminus Z
\]
that also identifies the intended metric structures. The metrics, their normalization, the precise choice of \(B\), and the required compatibility or extension of \(\Phi\) across singular loci are not specified in the supplied text. Thus this is a conditional statement of the reported target, not a complete formulation of the full conjecture.

\textbf{Scope and qualifications supplied with the problem.}
The excerpt reports the following background result as Theorem 2, citing its reference [11, Theorem 6.1]: for a minimal divisorial log terminal (dlt) model \(\mathcal X\) of \(X\), the retraction
\[
\rho_{\mathcal X}:X^{\mathrm{an}}\longrightarrow\mathrm{Sk}(X)
\]
is an \(n\)-dimensional affinoid torus fibration away from \(Z\). Here an affinoid torus fibration is the non-Archimedean analogue of a smooth torus fibration; its technical definition and the full hypotheses of the cited theorem are not included. This supplied background assertion is distinct from the requested metric comparison.

The excerpt then specializes to Calabi--Yau hypersurface degenerations
\[
X_t=\left\{z_0z_1\cdots z_{n+1}+t f_{n+2}(z)=0\right\}\subset\PP^{n+1},
\]
where \(f_{n+2}\) is a generic homogeneous polynomial of degree \(n+2\). The ground field and the meaning of genericity are not specified. For \(n=3\), these are quintic threefolds. The excerpt attributes to work in collaboration with L\'eonard Pille-Schneider, identified in the supplied literature assessment as Mazzon--Pille-Schneider, the construction of new non-Archimedean retractions giving evidence for the comparison conjecture. More precisely, its Theorem 3, citing [8, Theorem A], reports a continuous retraction
\[
\rho:X^{\mathrm{an}}\longrightarrow\mathrm{Sk}(X)
\]
which is an affinoid torus fibration outside a piecewise-linear locus \(\Gamma\) of codimension \(2\), contained in the \(2\)-skeleton of \(\mathrm{Sk}(X)\). The induced integral-affine structure on \(\mathrm{Sk}(X)\setminus\Gamma\) is reported to be isomorphic to structures constructed in references [1], [2], [6], and [12] of the original report. The excerpt describes the corresponding base \(B\) through dual complexes of toric degenerations in the Gross--Siebert program and through a Gromov--Hausdorff limit. It supplies no precise metric comparison theorem. The loci \(\Gamma\) and \(Z\) are not asserted here to coincide.

The reported construction glues retractions associated with several minimal dlt models adapted to different regions of the skeleton. The excerpt also states that the displayed hypersurface model \(\mathcal X\) has singularities that are not \(\Q\)-factorial, so its associated retraction is only well-defined over the union \(U\) of open maximal faces of \(\mathrm{Sk}(X)\). It reports that this suffices to prove a conjecture of Li relating SYZ fibrations to the non-Archimedean retraction over \(U\), without supplying that conjecture's statement. Separately, it cites [7, Theorem 1.3] for existence of SYZ fibrations on large open regions provided a comparison property holds for the solution of a non-Archimedean Monge--Amp\`ere equation over \(U\), referring also to [7, Proposition 3.7] and [14]. These conditional and regional statements do not supply a proof of the full metric integral-affine comparison.
\subsection{Short English statement}
The reported Kontsevich--Soibelman problem asks whether the regular part of a limiting SYZ base agrees with the regular part of the essential skeleton, preserving both integral-affine and metric structures. The excerpt reports integral-affine comparison results for quintic threefolds as evidence. Its omitted hypotheses, omitted Conjecture 1, and unspecified metrics prevent a fully precise universal formulation. The supplied partially solved label is dated curation, not an independently verified conclusion.
\subsection{Sources}
\begin{itemize}
\item Local UnsolvedMath registry, supplied record with permanent ID \(30005174\), source problem code \texttt{catalogue entry}, title ``Kontsevich--Soibelman Conjecture for Calabi--Yau Degenerations.'' The record belongs to the Oberwolfach Reports Open Problems collection and lists proposed year \(2022\); no proposer is supplied. Its registry status is \texttt{partially\_solved}. Its \texttt{corrected\_verified} statement label records supplied editorial verification, not independent mathematical review.
\item Supplied report citation: \emph{Algebraic Geometry: Moduli Spaces, Birational Geometry and Derived Aspects} (2023), Oberwolfach Reports, pp.~1805--1864, DOI \url{https://doi.org/10.4171/owr/2022/32}. The supplied excerpt carries page number 1833. The literature-source label is ``Oberwolfach Report 32/2022.''
\item Supplied literature source: Mazzon--Pille-Schneider, \emph{Toric geometry and integral affine structures}, \url{https://arxiv.org/abs/2110.04223}.
\item Literature assessment dated exactly \texttt{2026-08-22}: the supplied curation states that Mazzon--Pille-Schneider construct retractions and match integral-affine structures for quintic threefolds, providing evidence rather than a proof of the full Kontsevich--Soibelman conjecture. This assessment is supplied dated curation derived from the OWR record and has not been independently verified here.
\item Attribution and internal references retained from the excerpt: Kontsevich and Soibelman for the comparison conjecture; collaboration with L\'eonard Pille-Schneider for the quintic retractions; the Gross--Siebert program for the toric-degeneration description; and Li for the stated conditional connection with SYZ existence. Internal labels [11, Theorem 6.1], [8, Theorem A], [1], [2], [6], [12], [7, Theorem 1.3], [7, Proposition 3.7], and [14] refer to the report's bibliography, which is not supplied. Their bibliographic identities have not been reconstructed.
\end{itemize}
\end{document}
